\documentclass[a4paper]{article}
\usepackage{graphicx}
\usepackage{amsmath,amssymb}
\usepackage{hyperref}
\usepackage{minted}
\usepackage{multirow}

\title{}
\date{}
\begin{document}
    \maketitle
    It reminds me of a programming language called Lisp \textbf{}
    
    
    \section{Usecases}
        Decorator in python canz both be used for
        \begin{itemize}
            \item overwriting methods and properties similar to inheritence but from core to skin
            \item add extra methods and properties
        \end{itemize}
    
    
    \section{Downsides}
        Repetition of method definitions in Decorator
        
    \section{UML}
    
    \section{usecase:}
        \begin{itemize}
            \item Adding methods to a basic object
            \item Changing the behavior of a basic object
        \end{itemize}
    
    \section{Implementation}
        \begin{itemize}
            \item Interface: The genetic of what is trying to be modeled
            \item Abstract class \textbf{Decoratr} inherits(implements) from \textbf{Interface}, this class is the genetics of being the concept of a decoration. It composes a protected \textbf{\_inner} property of type Interface. What ever method or property of it is called, it calls the same method on the \textbf{\_inner}
            \item Decorations are the classes who inherit from \textit{Decorator}. Either they implement a method or from the Decorator or if that method doesnt exist automatically , according to inheritance rules, once again \textbf{Decorator.\_inner.method or .property} will send it to next, more inner method or property. Another option is
        \end{itemize}
    
    \section{Notes}
        \begin{itemize}
            \item Decoration reimplements inheritince in horizentally as real inheritance does it vertically. Or maybe it is better to say that it gives us a second utility for inhertitence
            \item Decorations are not needed to implement \textbf{\_\_init\_\_} unless they neeed their own arguments for the \textbf{\_\_init\_\_(self, inner:Decorator, a, b, c)} in that case \textbf{super().\_\_init(inner)\_\_} must be called in \textbf{\_\_init\_\_} body.
            \item In case a special kind of decorator can only decorate another special then \textbf{\_\_init\_\_(inner:SpecialDecoration, ...)}
            \item if two decorators can not be used together then you probably need two set of decorators. or you need several sub classes from the Basic class
            \item the docoration does not need to use \mintinline{python}{self.\_inner.method()}, it just calls \mintinline{python}{self.method()} and it calss
            \item  it is always better to used passive form names for decoratars for example \"diced\" for something that adds dictionary abilities to inner object
        \end{itemize}
    
    \section{When to add to \_inner, when to the self in every concerete decorator?}
        \textbf{“If I remove the decorator, should this thing still exist?”
        \\
        If yes → use \_inner
        \\
        If no → use self
        }
    
    \section{Downsides}
        \begin{itemize}
            \item sometime two decorators are inconsistenet with each other
        \end{itemize}
    
    \section{Ralation with natural language}
        \begin{itemize}
            \item The question is if and how we describe decorated concepts in natural language
        \end{itemize}
\end{document}